\documentclass[11pt]{article}
\usepackage{mathblatt}
\begin{document}
\weit
\typeout{PROBE-BEGIN koerper-e1-k1-s0-v1}
\begin{aufgabe}{\detokenize{koerper-e1-k1-s0-v1}}
\begin{teile}
\teil Probe

\prismadreieck{4}{2}{2.5}{}{}{}
\end{teile}
\end{aufgabe}
\typeout{PROBE-END koerper-e1-k1-s0-v1}
\typeout{PROBE-BEGIN koerper-e1-k1-s0-v2}
\begin{aufgabe}{\detokenize{koerper-e1-k1-s0-v2}}
\begin{teile}
\teil Probe

\prismadreieck{3}{3}{1.5}{}{}{}
\end{teile}
\end{aufgabe}
\typeout{PROBE-END koerper-e1-k1-s0-v2}
\typeout{PROBE-BEGIN koerper-e1-k1-s1-v2}
\begin{aufgabe}{\detokenize{koerper-e1-k1-s1-v2}}
\begin{teile}
\teil Probe

\prismadreieck{4}{2.5}{2}{}{}{}
\end{teile}
\end{aufgabe}
\typeout{PROBE-END koerper-e1-k1-s1-v2}
\typeout{PROBE-BEGIN koerper-e1-k5-s4-v2}
\begin{aufgabe}{\detokenize{koerper-e1-k5-s4-v2}}
\begin{teile}
\teil Probe

\prismadreieck{4}{2.5}{2}{}{}{} \rechenplatz[halb]{4}
\end{teile}
\end{aufgabe}
\typeout{PROBE-END koerper-e1-k5-s4-v2}
\typeout{PROBE-BEGIN koerper-e3-k1-s0-v2}
\begin{aufgabe}{\detokenize{koerper-e3-k1-s0-v2}}
\begin{teile}
\teil Probe

\prismadreieck{3}{2}{4}{}{}{}
\end{teile}
\end{aufgabe}
\typeout{PROBE-END koerper-e3-k1-s0-v2}
\end{document}
